\begin{abstract}
    Scientists often want to explain why an outcome is different in two
populations.
For instance, if patient outcomes are different across two neighborhoods, the outcome difference could be due to differences in the patients themselves (covariates) or differences in medical care access (outcomes given covariates).
%For instance, if patient outcomes are different at two hospitals, the different outcomes could be due to differences in how easy the patients are to treat (particular covariates) or differences in care at the two hospitals (outcomes given covariates).
%a jobs program benefits members of one city more than another, the benefits difference could be due to differences in program participants (particular covariates) or differences in the local
%labor markets (outcomes given covariates).
The Oaxaca--Blinder decomposition (OBD) is a standard tool in econometrics that explains the difference in the mean outcome
across two populations. Since the original Oaxaca--Blinder (OB) papers, it has been well known that the OBD is not symmetric with respect to the two populations; depending on the choice of a reference population, it is possible to derive numerically different answers. In the present paper, we provide what is, to the best of our knowledge, the first example illustrating that major conclusions can change between the two natural directions of the OBD. In a simulation inspired by real-life health data \mqedit{and real data?}, we illustrate a sign difference: one direction concludes that quality of healthcare improves patient outcomes, and the other concludes quality of care worsens patient outcomes. While model misspecification is a potential concern with the OBD, our use of a simulation \mqedit{and real?} demonstrates the problem can arise even under correct specification. Our supporting theory suggests that this sign change is possible in a wide variety of scenarios and not specific to our concrete example.
  \end{abstract}